\section{The Navier--Stokes equations}\label{sec:ns}

The second instantiation is the incompressible Navier--Stokes equations, with
$d = 3$ and $u_i$ the velocity field. The framework applies unchanged; what
this section supplies is the clock, the nondimensional right-hand side, and
the structure of the physics operator. Two items needed to actually run the
algorithm of \S\ref{sec:algorithm} --- the time-marching scheme inside
$\Phi^{(\mathrm{P})}$ and the constraint functionals --- remain open, and are
marked as such below.

\subsection{Governing equations and clock}\label{ssec:ns-clock}

The dimensional equations are
\begin{equation}\label{eq:ns-momentum}
  \frac{\partial u_i^*}{\partial t^*} + u_j^* \frac{\partial u_i^*}{\partial x_j^*}= -\frac{1}{\rho}\frac{\partial p^*}{\partial x_i^*} + \nu \frac{\partial^2 u_i^*}{\partial x_j^* \partial x_j^*},
\end{equation}
together with the continuity equation,
\begin{equation}\label{eq:ns-continuity}
  \frac{\partial u_i^*}{\partial x_i^*} = 0,
\end{equation}
so that, in the notation of \eqref{eq:fw-general-pde},
\begin{equation}\label{eq:ns-R}
  R_i^*(u^*) = - u_j^* \frac{\partial u_i^*}{\partial x_j^*}
  - \frac{1}{\rho}\frac{\partial p^*}{\partial x_i^*}
  + \nu \frac{\partial^2 u_i^*}{\partial x_j^* \partial x_j^*}.
\end{equation}
In addition to the rescaling \eqref{eq:fw-nondim} of $u_i$ and $x_j$, the
pressure is nondimensionalized by the dynamic pressure,
\begin{equation}\label{eq:ns-pressure-nondim}
  p = \frac{p^*}{\rho U^2}.
\end{equation}
Term by term, the convective and pressure terms scale as $U^2/L$,
\begin{equation}\label{eq:ns-convection}
  u_j^* \frac{\partial u_i^*}{\partial x_j^*} = \frac{U^2}{L} u_j \frac{\partial u_i}{\partial x_j},
  \qquad
  \frac{1}{\rho}\frac{\partial p^*}{\partial x_i^*} = \frac{U^2}{L} \frac{\partial p}{\partial x_i},
\end{equation}
while the viscous term scales as $\nu U / L^2$,
\begin{equation}\label{eq:ns-viscous}
  \nu \frac{\partial^2 u_i^*}{\partial x_j^* \partial x_j^*} = \frac{\nu U}{L^2}\frac{\partial^2 u_i}{\partial x_j \partial x_j}.
\end{equation}
The clock is chosen so that the dominant, advective part of the physical term
of the lemma \eqref{eq:fw-lemma} carries unit coefficient,
\begin{equation}\label{eq:ns-clock}
  \frac{\tau}{U} \cdot \frac{U^2}{L} = 1
  \quad\Longrightarrow\quad
  \tau = \frac{L}{U},
\end{equation}
which is the eddy turnover time of the current scales. In log form, as
required by \eqref{eq:fw-ln-dtstar},
\begin{equation}\label{eq:ns-ln-clock}
  \ln \tau_k = \ln L_k - \ln U_k .
\end{equation}
Applying $\tau / U = L/U^2$ to \eqref{eq:ns-R} and introducing the
instantaneous Reynolds number
\begin{equation}\label{eq:ns-reynolds}
  Re = \frac{U L}{\nu},
\end{equation}
the rescaling lemma \eqref{eq:fw-lemma} gives the nondimensional momentum
equation
\begin{equation}\label{eq:ns-nondim-momentum-base}
  \frac{\partial u_i}{\partial t}
  = \sigma_L\, x_j \frac{\partial u_i}{\partial x_j}
  - \sigma_U\, u_i
  - u_j \frac{\partial u_i}{\partial x_j}
  - \frac{\partial p}{\partial x_i}
  + \frac{1}{Re} \frac{\partial^2 u_i}{\partial x_j \partial x_j},
\end{equation}
and the continuity equation is unchanged by the rescaling,
\begin{equation}\label{eq:ns-nondim-continuity}
  \frac{\partial u_i}{\partial x_i} = 0.
\end{equation}

\subsection{Viscosity}\label{ssec:ns-viscosity}
Let the kinematic viscosity be given by
\begin{equation}\label{eq:ns-viscosity-split}
  \nu = \nu_{\mathrm{turb}} + \nu_{\mathrm{mol}},
\end{equation}
where $\nu_{\mathrm{mol}}$ is the molecular viscosity and $\nu_{\mathrm{turb}}$
is a turbulent viscosity representing under-resolved eddies, approximated as
\begin{equation}\label{eq:ns-turb-viscosity}
  \nu_{\mathrm{turb}} = \Delta \cdot U L,
\end{equation}
where $\Delta$ is a parameter related to the resolution of the mesh. In the
limit of infinite Reynolds number, where $\nu_{\mathrm{mol}} = 0$, we have
$1/Re = \Delta$ and the momentum equation
\eqref{eq:ns-nondim-momentum-base} becomes
\begin{equation}\label{eq:ns-nondim-momentum}
  \frac{\partial u_i}{\partial t}
  = \sigma_L\, x_j \frac{\partial u_i}{\partial x_j}
  - \sigma_U\, u_i
  - u_j \frac{\partial u_i}{\partial x_j}
  - \frac{\partial p}{\partial x_i}
  + \Delta \cdot \frac{\partial^2 u_i}{\partial x_j \partial x_j}.
\end{equation}

\subsection{Fourier representation}\label{ssec:ns-fourier}
Using the Fourier convention \eqref{eq:fw-fourier-convention}, the viscous
term becomes
\begin{equation}\label{eq:ns-fourier-viscous}
  \widehat{\frac{\partial^2 u_i}{\partial x_j \partial x_j}} = -k^2 \hat{u}_i,
\end{equation}
and the pressure gradient becomes
\begin{equation}\label{eq:ns-fourier-pressure}
  \widehat{\frac{\partial p}{\partial x_i}} = \mathrm{i} k_i \hat{p}.
\end{equation}
The convective term
\begin{equation}\label{eq:ns-fourier-convection}
  \hat{N}_i = \widehat{u_j \frac{\partial u_i}{\partial x_j}},
\end{equation}
will be evaluated pseudo-spectrally. The scaling terms transform as in
\eqref{eq:fw-scaling-pde-fourier} with $d = 3$, so the transform of
\eqref{eq:ns-nondim-momentum} is
\begin{equation}\label{eq:ns-nondim-momentum-fourier}
  \frac{\partial \hat{u}_i}{\partial t}
  = -\sigma_L \left( k_j \frac{\partial \hat{u}_i}{\partial k_j} + 3\, \hat{u}_i \right)
  - \sigma_U\, \hat{u}_i
  - \hat{N}_i
  - \mathrm{i} k_i \hat{p}
  - \Delta \cdot k^2 \hat{u}_i,
\end{equation}
and the transform of the continuity equation
\eqref{eq:ns-nondim-continuity} is
\begin{equation}\label{eq:ns-fourier-continuity}
  k_i \hat{u}_i = 0.
\end{equation}

\subsection{Pressure and the solenoidal projection}\label{ssec:ns-projection}
Contracting \eqref{eq:ns-nondim-momentum-fourier} with $k_i$ and using the
continuity constraint gives
\begin{equation}\label{eq:ns-pressure-contracted}
  0 = -\sigma_L \left( k_i k_j \frac{\partial \hat{u}_i}{\partial k_j} + d\, k_i \hat{u}_i \right)
  - \sigma_U\, k_i \hat{u}_i
  - k_i \hat{N}_i
  - \mathrm{i} k^2 \hat{p}
  - \Delta \cdot k^2 k_i \hat{u}_i.
\end{equation}
The remaining scaling contribution also vanishes: differentiating
\eqref{eq:ns-fourier-continuity} with respect to $k_j$ gives
\begin{equation}\label{eq:ns-continuity-dk}
  \frac{\partial}{\partial k_j}\left( k_i \hat{u}_i \right)
  = \hat{u}_j + k_i \frac{\partial \hat{u}_i}{\partial k_j} = 0
  \quad\Longrightarrow\quad
  k_i \frac{\partial \hat{u}_i}{\partial k_j} = -\hat{u}_j,
\end{equation}
so that $k_i k_j \partial \hat{u}_i / \partial k_j = -k_j \hat{u}_j = 0$.
Hence the scaling terms are divergence free on their own --- the self-similar
frame does not generate pressure --- and \eqref{eq:ns-pressure-contracted}
reduces to a Poisson equation for the pressure,
\begin{equation}\label{eq:ns-fourier-poisson}
  \mathrm{i} k^2 \hat{p} = -k_i \hat{N}_i
  \quad\Longrightarrow\quad
  \hat{p} = \mathrm{i}\,\frac{k_l \hat{N}_l}{k^2},
  \qquad k \neq 0,
\end{equation}
with $\hat{p}$ at $k = 0$ fixed arbitrarily, since only its gradient enters
the momentum equation. Substituting \eqref{eq:ns-fourier-poisson} into
\eqref{eq:ns-nondim-momentum-fourier}, the pressure term becomes
\begin{equation}\label{eq:ns-fourier-pressure-substituted}
  -\mathrm{i} k_i \hat{p} = \frac{k_i k_l}{k^2} \hat{N}_l,
\end{equation}
which combines with $-\hat{N}_i$ to give the solenoidal projection of the
convective term. Defining the projection tensor
\begin{equation}\label{eq:ns-projection-tensor}
  P_{il}(\mathbf{k}) = \delta_{il} - \frac{k_i k_l}{k^2},
\end{equation}
the evolution equation splits exactly along the lines of the lemma
\eqref{eq:fw-lemma},
\begin{equation}\label{eq:ns-fourier-final}
  \frac{\partial \hat{u}_i}{\partial t}
  = \underbrace{-\sigma_L \left( k_j \frac{\partial \hat{u}_i}{\partial k_j} + 3\, \hat{u}_i \right)
  - \sigma_U\, \hat{u}_i}_{\text{scaling terms}}
  \underbrace{- P_{il}(\mathbf{k})\, \hat{N}_l
  - \Delta \cdot k^2 \hat{u}_i}_{\text{physical terms}}.
\end{equation}

\subsection{The physics step}\label{ssec:ns-physics-step}

The physics operator \eqref{eq:fw-physics-flow} is the flow of the bracketed
physical terms of \eqref{eq:ns-fourier-final},
\begin{equation}\label{eq:ns-physics-flow}
  \Phi_{\Delta t}^{(\mathrm{P})}\bigl(\hat{u}_i(t)\bigr)
  = \hat{u}_i(t) - \int_{t}^{t+\Delta t}
    \left[ P_{il}(\mathbf{k})\, \hat{N}_l + \Delta \cdot k^2 \hat{u}_i \right] \mathrm{d}t .
\end{equation}

\emph{Open.} Unlike the 1D case, \eqref{eq:ns-physics-flow} has no closed form:
the convective term is nonlinear, so $\Phi^{(\mathrm{P})}$ must be realized by
an actual time-marching scheme. That scheme has not been chosen. Whatever is
chosen must meet the accuracy requirement of \S\ref{ssec:alg-substeps} ---
second order for the Strang composition to retain its order --- and must
handle the dealiasing of the pseudo-spectral product
\eqref{eq:ns-fourier-convection}. The linear term $-\Delta k^2 \hat u_i$ admits
an exact integrating factor and need not be marched with the rest.

\subsection{Constraint functionals}\label{ssec:ns-constraints}

\emph{Open.} The constraint functionals of \S\ref{ssec:fw-constraints} have not
been chosen for the Navier--Stokes case. The moments
\eqref{eq:1d-functionals} used for the 1D equation do not carry over: a
velocity field in a periodic box has $\int u_i\, \mathrm{d}\mathbf{x} \approx 0$,
so $F_1$ as written is not usable as a scale. Any admissible replacement must
be a pair $(F_1, F_2)$ that is homogeneous under the scaling
\eqref{eq:fw-scale-step} in the sense of \eqref{eq:fw-homogeneity}, with
exponents satisfying the solvability condition
\eqref{eq:fw-constraint-solvable} --- that is, one functional sensitive
primarily to the amplitude of the field and one sensitive primarily to its
length scale. Quantities of the right character are the kinetic energy and an
integral length scale extracted from the energy spectrum, but neither their
exponents nor their behaviour under the constrained step has been worked out
here.

\subsection{Boundary treatment}\label{ssec:ns-boundary}

The domain is a periodic box, so no boundary condition is imposed beyond
periodicity itself and the substep \textsc{ApplyBC} of
\S\ref{ssec:alg-substeps} is the identity. The tension noted in
\S\ref{ssec:1d-boundary} between a spectral discretization and a clamped edge
therefore does not arise here.
